%! Algebra 2 — Unit 1, Lesson 1 — Group Activity (Differentiated)
\documentclass[10pt]{article}
\usepackage{algebra2-article}
\usepackage{algebra2-boxes}

\newcommand{\TallMath}[1]{$\displaystyle #1\rule[-1.4em]{0pt}{3.2em}$}

%=============================================================================
\begin{document}

\pageheader{Unit 1, Lesson 1}{Group Activity}
\namepartnerperiod

\vspace{0.15in}

% ─────────────────────────────────────────────
%  TIER R — Remediate
% ─────────────────────────────────────────────
\begin{tcolorbox}[colback=white,colframe=black!40,boxrule=0.7pt,arc=2mm,
    title={\bfseries Tier R — Remediate},fonttitle=\small,breakable]

\begin{enumerate}

    \item Evaluate $2x + 3y$ when $x = 4$ and $y = -2$.

    \vspace{0.7in}

    \item Simplify. Write your answer as a single power of $x$.
    \[x^4 \cdot x^3\]

    \vspace{0.5in}

    \item Find the fully factored form of $x^2 + 5x + 6$.
    \quad \textit{Hint: what two numbers add to $5$ and multiply to $6$?}

    \vspace{0.7in}

    \item Which expression is the fully simplified form of $\sqrt{48}$?

    \begin{enumerate}[label=\Alph*.]
        \item $4\sqrt{3}$
        \item $2\sqrt{12}$
        \item $\sqrt{16 \cdot 3}$
        \item $6\sqrt{2}$
    \end{enumerate}

\end{enumerate}
\end{tcolorbox}

\vspace{0.2in}

% ─────────────────────────────────────────────
%  TIER A — Approaching Proficiency
% ─────────────────────────────────────────────
\begin{tcolorbox}[colback=white,colframe=black!40,boxrule=0.7pt,arc=2mm,
    title={\bfseries Tier A — Approaching Proficiency},fonttitle=\small,breakable]

\begin{enumerate}

    \item Evaluate $3x^2 - 2y + 4z$ when $x = -2$, $y = 5$, and $z = -3$.

    \vspace{0.8in}

    \item Which expression is the equivalent, fully simplified form of
    \TallMath{\dfrac{x^3 \cdot x^5}{x^2}}?

    \begin{enumerate}[label=\Alph*.]
        \item $x^6$
        \item $x^8$
        \item \TallMath{\dfrac{x^8}{x^2}}
        \item $x^{15}$
    \end{enumerate}

    \vspace{0.1in}

    \item Find the fully factored form of each expression.

    \begin{enumerate}[label=\alph*.]
        \item $x^2 + 9x + 20$

        \vspace{0.75in}

        \item $6x^2 + 11x + 3$ \quad \textit{(use slip-and-slide)}

        \vspace{0.75in}
    \end{enumerate}

    \item Which expression is the equivalent, fully simplified form of $\sqrt{180}$?

    \begin{enumerate}[label=\Alph*.]
        \item $6\sqrt{5}$
        \item $3\sqrt{20}$
        \item $\sqrt{36 \cdot 5}$
        \item $2\sqrt{45}$
    \end{enumerate}

\end{enumerate}
\end{tcolorbox}

% ─────────────────────────────────────────────
%  TIER E — Extension
% ─────────────────────────────────────────────
\begin{tcolorbox}[colback=white,colframe=black!40,boxrule=0.7pt,arc=2mm,
    title={\bfseries Tier E — Extension},fonttitle=\small,breakable]

\begin{enumerate}

    \item Evaluate \TallMath{\dfrac{2a^3 - b}{a^2 + b^2}} when $a = -3$ and $b = 4$.

    \vspace{0.9in}

    \item Simplify completely. Write your answer as a single power of $x$.
    \TallMath{\dfrac{(2x^3)^2 \cdot x^{-1}}{4x^5}}

    \vspace{0.9in}

    \item Factor each expression completely.

    \begin{enumerate}[label=\alph*.]
        \item $4x^2 - 25$

        \vspace{0.75in}

        \item $10x^2 - 13x - 3$

        \vspace{0.75in}
    \end{enumerate}

    \item Simplify $\sqrt{75} - \sqrt{12}$. Show each step.

    \vspace{0.9in}

\end{enumerate}
\end{tcolorbox}

\end{document}
